\section*{Chapter \ref{ch:b1:open-economy} --- The Open Economy}
% Numbers recomputed in code/bachelor-1/27-open-economy/exercises/solutions.py (tested).

\begin{solution}{exo:b1:open-economy:nx}
$NX = S - I = 50$.
\end{solution}

\begin{solution}{exo:b1:open-economy:real-ex}
$\varepsilon_r = 1.25 \times 100/110 = 1.136$: a unit of Kestria's goods costs 1.136 units of Brenland's.
\end{solution}

\begin{solution}{exo:b1:open-economy:read}
Netherlands $0.46$, France $2.76$, Italy $5.86$. The Dutch peg was believed: its
rate barely exceeded Germany's. Markets expected the lira, and to a lesser
degree the franc, to be devalued from time to time, and required a premium.
\end{solution}

\begin{solution}{exo:b1:open-economy:predict}
Saving falls, so $NX = S - I(r^*)$ falls, and the real exchange rate falls: a real
appreciation.
\end{solution}

\begin{solution}{exo:b1:open-economy:small}
$I = 500 - 80 = 420$; $NX = 400 - 420 = -20$; $50\varepsilon_r - 100 = -20$ gives $\varepsilon_r = 1.6$.
\end{solution}

\begin{solution}{exo:b1:open-economy:float}
Output is fixed by LM$^*$: no change. Net exports must fall by 50, so $e$ falls by
$50/25 = 2$: the currency appreciates and crowds out net exports one for one.
\end{solution}

\begin{solution}{exo:b1:open-economy:peg}
The purchase raises $M$ and pushes $r$ below $r^*$; capital flows out, and holders of
home currency sell it for foreign currency. To hold $e$, the central bank buys
its own currency with reserves, which lowers $M$. This goes on until $M$ is back
where it was: reserves are lower, nothing else has changed.
\end{solution}

\begin{solution}{exo:b1:open-economy:ml}
$0.6 + 0.7 = 1.3 > 1$: yes. $0.3 + 0.4 = 0.7 < 1$: no; the higher price of imports
outweighs the change in volumes, and net exports fall.
\end{solution}

\begin{solution}{exo:b1:open-economy:count}
Under a credible peg with free capital flows, the gap should be near zero.
Gaps of several points over many years mean the pegs were not fully
believed --- both countries devalued more than once --- or that capital
was not free to move; the largest gaps came when devaluations were
expected.
\end{solution}

\begin{solution}{exo:b1:open-economy:derive}
Differentiate: $(1 - C')dY = dG + NX'de$ and $L_Y dY = d(M/P)$. Float ($M$ given): $dY = d(M/P)/L_Y$, and
for $dG$ alone $dY = 0$, $de = -dG/NX'$. Peg ($de = 0$): $dY = dG/(1 - C')$, and the money supply must
change by $d(M/P) = L_Y dY$; an independent change in $M$ is undone by capital flows.
\end{solution}

\begin{solution}{exo:b1:open-economy:sweden}
Free capital flows equalise expected returns: holding home assets earns $r$,
holding foreign assets earns $r^*$ plus the expected rise in $e$; so $r = r^* + \Delta e^e$.
$10\%$ a week is $52 \times 10 = 520\%$ a year at simple interest, plus $r^*$: the order of the
Swedish rate of September 1992.
\end{solution}

\begin{solution}{exo:b1:open-economy:refute}
A country with a credible fixed rate and open capital markets whose
interest rate moved independently of the base country's, persistently and
by a large margin. The Dutch rate shows the opposite.
\end{solution}

\begin{solution}{pb:b1:open-economy:peg}
\textbf{1.} $129{,}777 - 128{,}664 = 1{,}113$ million dalens.

\textbf{2.} $393{,}264 - 225{,}014 - 78{,}653 = 89{,}597$.

\textbf{3.} $89{,}597 - 88{,}484 = 1{,}113$.

\textbf{4.} $157{,}306 - 165{,}171 = -7{,}865$.

\textbf{5.} $89{,}597 + 7{,}865 = 97{,}462$.

\textbf{6.} $40{,}231 - 34{,}739 = 5{,}492$ with Kestria; $1{,}113 - 5{,}492 = -4{,}379$ with the rest of the world.

\textbf{7.} $1.242/1.207 - 1 = 2.9\%$.

\textbf{8.} $I = 103.484 - 15 = 88.484$; $NX = 89.597 - 88.484 = 1.113$; $\varepsilon_r = (1.113 + 48.887)/50 = 1$.

\textbf{9.} $I = 83.484$; $NX = 6.113$.

\textbf{10.} $\varepsilon_r = 55/50 = 1.1$: a real depreciation of $10\%$.

\textbf{11.} Saving falls to 79.597: $NX = -8.887$ and $\varepsilon_r = 40/50 = 0.8$.

\textbf{12.} Output is at its natural level, set by capital, labour and
technology.

\textbf{13.} Saving and investment move together far more than $r = r^*$ with
perfect mobility allows: the coefficient was $0.89$ in 1960--1974 and is still
$0.272$ in 2000--2019 for 18 of the 21 countries of their table.

\textbf{14.} With $c_0 = 37.056$: IS$^*$ $Y = 2.5(107.306 + 50e) = 268.265 + 125e$; LM$^*$ $181.632 = 0.5Y - 15$, so
$Y = 393.264$, and $e = 1$.

\textbf{15.} $Y = 393.264$; IS$^*$ shifts right by 25, so $e = 0.8$.

\textbf{16.} $Y = 2(191.632 + 15) = 413.264$; $e = (413.264 - 268.265)/125 = 1.16$.

\textbf{17.} $Y = 393.264 + 2.5 \times 10 = 418.264$.

\textbf{18.} By $0.5 \times 25 = 12.5$: the central bank buys kests with new dalens to stop
the dalen appreciating.

\textbf{19.} Nothing in the end: capital flows out, the bank sells 10 of
reserves to hold the rate, and the money supply returns.

\textbf{20.} Monetary policy under a float, fiscal policy under a peg.

\textbf{21.} The kest return plus $5\%$ over the month.

\textbf{22.} $1.05^{12} - 1 = 79.6\%$.

\textbf{23.} $3 + 79.6 = 82.6\%$, about $83\%$, with the linear parity condition
$r = r^* + \Delta e^e$ (compounding both, $1.03 \times 1.05^{12} - 1 = 85.0\%$).

\textbf{24.} $I = 103.484 - 5 \times 82.6 = -309.4$: impossible; in the toy model investment would
collapse long before. The rate can be held for days, not months.

\textbf{25.} $4/1.242 = 3.22$ billion kests a day: $30/3.22 = 9.3$ days.

\textbf{26.} Holding the peg against a $5\%$ expected devaluation in one month
needs an annualised rate of about $83\%$: Kestria's $3\%$ plus a premium of 80
points ($79.6$).
\end{solution}
